% !PRINTED 522
\weberpage{522}

Hier stehen nun rechter und linker Hand ganze Functionen von $t$, deren Co\"efficienten ganze Functionale sind. Auf der rechten Seite haben alle diese Co\"efficienten den Factor $\mu$, also auch den Factor $\pi$, w\"ahrend nach Voraussetzung nicht alle Co\"efficienten von $\chi$ diesen Factor haben. Da $E$ eine Einheit ist, so enthalten auch die Co\"efficienten von $E$ den Factor $\pi$ nicht, und folglich k\"onnen nach dem Satze 1. auch die Co\"efficienten von $E\chi^m$ nicht alle durch $\pi$ theilbar sein, was doch die Gleichung (4) verlangen w\"urde. Daraus ergiebt sich, dass unsere Annahme, $\varphi$ sei ganz, $\varphi_0, \varphi_1, \ldots, \varphi_h$ dagegen nicht alle ganz, unstatthaft ist, und der Satz 2. ist somit bewiesen.

Nehmen wir nun an, in $\varphi$ seien die Co\"efficienten $\varphi_0, \varphi_1, \ldots$ selbst wieder ganze Functionen einer Variablen, und wenden den Satz 2. wiederholt darauf an, so gelangen wir zu dem Schlusse:

3. \gesperrt{Eine ganze rationale Function beliebig vieler Ver\"anderlicher, deren Co\"efficienten Zahlen oder Functionale mit anderen Variablen sind, ist nur dann ein ganzes Functional, wenn die Co\"efficienten ganz sind.}

Wir k\"onnen jedes Functional $\omega$ als Quotienten zweier ganzer Functionen in der Weise darstellen, dass der Nenner eine primitive Function im K\"orper $R$ wird.

Denn sind $\varphi, \psi$ ganze Functionen in $\Omega$, und ist

\[
\omega = \frac{\varphi}{\psi}, \quad N(\psi) = \psi\psi',
\]

so k\"onnen wir den Bruch $\omega$ durch $\psi'$ erweitern und erhalten im Nenner eine ganze Function mit rationalen Co\"efficienten. Den Theiler dieser Function k\"onnen wir dann zum Z\"ahler von $\omega$ rechnen, und erhalten, wenn $E$ eine primitive Function, $\chi$ eine ganze Function in $\Omega$ bedeutet:

\[
E\omega = \chi,
\]

d. h. man kann jedes Functional $\omega$ in $\Omega$ durch Multiplication mit einer primitiven Function $E$ in eine ganze Function der Variablen verwandeln, deren Co\"efficienten Zahlen in $\Omega$ sind.

Mit Zuziehung des Satzes 3. ergiebt sich hieraus:

4. \gesperrt{Jedes ganze Functional $\omega$ im K\"orper $\Omega$ ist associirt mit einer ganzen Function $\chi$, deren Co\"effi-}
%CONT>>
